\begin{theorembox}{Theorem}
\begin{theorem}[Hyperbolic Pythagorean identity]
\label{thm:hyperbolic-pythagorean-identity}
For every $t\in\mathbb{R}$,
\[
  \cosh^2 t-\sinh^2 t=1.
\]
\end{theorem}
\end{theorembox}
\begin{remark*}[Proof]
This is \hyperref[thm:fundamental-hyperbolic-identity]{the fundamental
hyperbolic identity}.
\end{remark*}
\begin{dependencies}
\begin{itemize}
  \item \hyperref[thm:fundamental-hyperbolic-identity]{Fundamental hyperbolic identity}
\end{itemize}
\end{dependencies}

\begin{theorembox}{Theorem}
\begin{theorem}[Exponential form of hyperbolic sine and cosine]
\label{thm:hyperbolic-exponential-form}
For every $t\in\mathbb{R}$,
\[
  \cosh t=\frac{e^t+e^{-t}}{2},
  \qquad
  \sinh t=\frac{e^t-e^{-t}}{2}.
\]
\end{theorem}
\end{theorembox}
\begin{remark*}[Proof]
Let $u=\cosh t$ and $v=\sinh t$. The identity $u^2-v^2=1$ gives
\[
  (u-v)(u+v)=1.
\]
On the right branch, $u+v>0$, so there is a unique real number $s$ such that
$u+v=e^s$. Then $u-v=e^{-s}$, and solving the two linear equations gives
\[
  u=\frac{e^s+e^{-s}}{2},
  \qquad
  v=\frac{e^s-e^{-s}}{2}.
\]
For the exponential parametrization
\[
  Q_s=\left(\frac{e^s+e^{-s}}{2},\frac{e^s-e^{-s}}{2}\right),
\]
a direct sector-area computation gives
$2\operatorname{HArea}(O,V,Q_s)=s$. Therefore the defining sector parameter is
$t=s$, and the displayed formulas follow. The negative case agrees with the
even/odd extension.
\end{remark*}
\begin{remark*}[Construction first]
This theorem is derived from the geometric construction. It is not the
definition of $\cosh$ and $\sinh$ in this chapter.
\end{remark*}
\begin{dependencies}
\begin{itemize}
  \item \hyperref[def:hyperbolic-sine]{Hyperbolic sine}
  \item \hyperref[def:hyperbolic-cosine]{Hyperbolic cosine}
  \item \hyperref[thm:fundamental-hyperbolic-identity]{Fundamental hyperbolic identity}
\end{itemize}
\end{dependencies}

\begin{theorembox}{Theorem}
\begin{theorem}[Hyperbolic sum formulas]
\label{thm:hyperbolic-sum-formulas}
For all $a,b\in\mathbb{R}$,
\[
  \cosh(a+b)=\cosh a\cosh b+\sinh a\sinh b,
\]
\[
  \sinh(a+b)=\sinh a\cosh b+\cosh a\sinh b.
\]
\end{theorem}
\end{theorembox}
\begin{remark*}[Proof]
Using the exponential form,
\[
  \cosh(a+b)
  =\frac{e^{a+b}+e^{-(a+b)}}{2}.
\]
Expanding the right side of the claimed identity gives
\[
\frac{e^a+e^{-a}}{2}\frac{e^b+e^{-b}}{2}
+\frac{e^a-e^{-a}}{2}\frac{e^b-e^{-b}}{2}
=\frac{e^{a+b}+e^{-(a+b)}}{2}.
\]
The proof for $\sinh(a+b)$ is the same expansion with the signs arranged so
that the mixed exponential terms survive with the claimed signs.
\end{remark*}
\begin{dependencies}
\begin{itemize}
  \item \hyperref[thm:hyperbolic-exponential-form]{Exponential form of hyperbolic sine and cosine}
\end{itemize}
\end{dependencies}

\begin{theorembox}{Theorem}
\begin{theorem}[Hyperbolic double-angle formulas]
\label{thm:hyperbolic-double-angle-formulas}
For every $t\in\mathbb{R}$,
\[
  \sinh(2t)=2\sinh t\cosh t,
\]
\[
  \cosh(2t)=\cosh^2 t+\sinh^2 t
  =2\cosh^2 t-1
  =1+2\sinh^2 t.
\]
\end{theorem}
\end{theorembox}
\begin{remark*}[Proof]
Set $a=b=t$ in the hyperbolic sum formulas. The first two displayed formulas
follow immediately. The remaining two forms of $\cosh(2t)$ follow by replacing
$\sinh^2 t$ with $\cosh^2 t-1$, or replacing $\cosh^2 t$ with
$1+\sinh^2 t$, using the hyperbolic Pythagorean identity.
\end{remark*}
\begin{dependencies}
\begin{itemize}
  \item \hyperref[thm:hyperbolic-sum-formulas]{Hyperbolic sum formulas}
  \item \hyperref[thm:hyperbolic-pythagorean-identity]{Hyperbolic Pythagorean identity}
\end{itemize}
\end{dependencies}

\begin{theorembox}{Theorem}
\begin{theorem}[Hyperbolic tangent sum formula]
\label{thm:hyperbolic-tangent-sum-formula}
For all $a,b\in\mathbb{R}$,
\[
  \tanh(a+b)=
  \frac{\tanh a+\tanh b}{1+\tanh a\tanh b}.
\]
\end{theorem}
\end{theorembox}
\begin{remark*}[Proof]
By definition,
\[
  \tanh(a+b)=\frac{\sinh(a+b)}{\cosh(a+b)}.
\]
Substitute the sum formulas:
\[
  \frac{\sinh a\cosh b+\cosh a\sinh b}
       {\cosh a\cosh b+\sinh a\sinh b}.
\]
Since $\cosh a$ and $\cosh b$ are nonzero, divide numerator and denominator by
$\cosh a\cosh b$ to obtain the displayed identity.
\end{remark*}
\begin{dependencies}
\begin{itemize}
  \item \hyperref[def:hyperbolic-tangent]{Hyperbolic tangent}
  \item \hyperref[thm:hyperbolic-sum-formulas]{Hyperbolic sum formulas}
\end{itemize}
\end{dependencies}
